\documentclass{template}

\usepackage{mathtools}
\pgfplotsset{compat=1.18}

\newcommand{\qedwhite}{\hfill \ensuremath{\Box}}


\begin{document}

\begin{center}
    \Large \textcolor{blue!70!black}{\textbf{Commutative Algebra Problem Set \#3}} \\
    \normalsize \textit{Dmytro Zakharov}
\end{center}

\begin{problem}[Exercise 1]
    Let $A=\mathbb{Z}$ and $M = \mathbb{Q} / \mathbb{Z}$. Show that 
    $M(x) = M \otimes_A \kappa(x)$ is zero for every $x \in \text{Spec}(A)$.
\end{problem}

\textbf{Solution.} First, consider $x=0$ separately:
\begin{equation*}
    M(0) = M \otimes_{\mathbb{Z}} \kappa(0) = M \otimes_{\mathbb{Z}} \text{Frac}(\mathbb{Z})
    \cong (\mathbb{Q} / \mathbb{Z}) \otimes_{\mathbb{Z}} \mathbb{Q} \cong 0
\end{equation*}

The last isomorphism comes from the fact that for any $\frac{n}{m}+\mathbb{Z} \in \mathbb{Q}/\mathbb{Z}$
and $\frac{p}{q} \in \mathbb{Q}$:
\begin{equation*}
    \left(\frac{n}{m} + \mathbb{Z}\right) \otimes \frac{p}{q} = \left(\frac{n}{m} + \mathbb{Z}\right) \otimes m \cdot \frac{p}{qm} = m\left(\frac{n}{m} + \mathbb{Z}\right) \otimes \frac{p}{qm} = 0.
\end{equation*} 

Now consider any non-zero prime ideal $p\mathbb{Z}$ corresponding to $x$. Then, 
$\kappa(x) = \text{Frac}(A/p\mathbb{Z}) = \text{Frac}(\mathbb{Z}/p\mathbb{Z}) \cong \text{Frac}(\mathbb{F}_p) \cong \mathbb{F}_p$. 
Thus we consider:
\begin{equation*}
    M(x) = (\mathbb{Q}/\mathbb{Z}) \otimes_{\mathbb{Z}} \mathbb{F}_p.
\end{equation*}

This is again clearly $0$. For arbitrary $\frac{n}{m}+\mathbb{Z} \in
\mathbb{Q}/\mathbb{Z}$ and $\alpha \in \mathbb{F}_p$ we have:
\begin{equation*}
    \left(\frac{n}{m}+\mathbb{Z}\right) \otimes \alpha = p\left(\frac{n}{mp}+\mathbb{Z}\right) \otimes \alpha = \left(\frac{n}{mp}+\mathbb{Z}\right) \otimes p\alpha = 0.
\end{equation*}

\begin{problem}[Exercise 2]
    Let $A$ be a ring and $M$ a finitely presented $A$-module: that is, exists a 
    presentation
        \[
    \begin{tikzcd}
        A^m \arrow[-Stealth]{r}{\alpha} & A^n \arrow[-Stealth]{r}{\pi} & M \arrow[-Stealth]{r}{} & 0,
    \end{tikzcd}
        \]
    or equivalently there exists a surjection $A^n \to M$ whose kernel is finitely 
    generated. Show that the kernel of \textbf{every} surjection $A^k \to M$ is 
    finitely generated.
\end{problem}

\textbf{Proof.} As noted in the hint, we first show that if
\begin{equation*}
    0 \to M' \xrightarrow{\iota} M \xrightarrow{\pi} M'' \to 0
\end{equation*}

is exact and both $M'$ and $M''$ are finitely generated, then so is $M$. Let
$M'$ be generated by $x_1,\dots,x_n$ while $M''$ by $y_1,\dots,y_m$. Let
$z_1,\dots,z_m$ be pre-images of $y_1,\dots,y_m$ in $M$ (that is,
$\pi(y_i)=z_i$). We claim $M = \langle x_1,\dots,x_n,z_1,\dots,z_m \rangle$, and
so $M$ is finitely generated. 

Indeed, take any $v \in M$. Note that $\pi(v) = \sum_{i=1}^m \alpha_i y_i$ for
$\alpha_i \in A$ since $M''$ is finitely generated. Thus trivially, $v -
\sum_{i=1}^m \alpha_i z_i \in \text{ker}(\pi) \cong \text{im}(\iota)$, and 
so we can write this difference using generators of $M'$:
\begin{equation*}
    v - \sum_{i=1}^m \alpha_i z_i = \sum_{j=1}^n \beta_i x_i, \quad \beta_i \in A
\end{equation*}

Thus, we have found the linear combination $\sum_{i=1}^m\alpha_i z_i +
\sum_{i=1}^n \beta_i x_i$ for any $v \in M$.

Now, we shall show that \textbf{every} surjection $A^k \to M$ is finitely
generated. Build the following commutative diagram:
\[
    \begin{tikzcd}
        & A^m \arrow[-Stealth]{d}{g} \arrow[-Stealth]{r}{} & A^n \arrow[-Stealth]{d}{f} \arrow[-Stealth]{r}{} & M \arrow[equal]{d}{} \arrow[-Stealth]{r}{} & 0 \\
        0 \arrow[-Stealth]{r}{} & K \arrow[-Stealth]{r}{} & A^k \arrow[-Stealth]{r}{\pi} & M \arrow[-Stealth]{r}{} & 0
    \end{tikzcd}
\]

Here, $f: A^n \to A^k$ is constructed by lifting the identity map $\text{id}_M$,
and the map $g: A^m \to K$ is obtained by restricting the map $f$. Now apply 
the Snake's lemma to get exact sequence:
\begin{equation*}
    0 \to \text{coker}(g) \to \text{coker}(f) \to 0.
\end{equation*}

Thus, $\text{coker}(f) \cong \text{coker}(g)$, and so we have the exact sequence:
\begin{equation*}
    0 \to \text{im}(g) \to K \to \text{coker}(f) \to 0,
\end{equation*} 

where both $\text{im}(g)$ and $\text{cokeg}(f)$ are finitely generated. Thus,
$K$ is finitely generated as well. \qedwhite

\begin{problem}[Exercise 3]
    Let $A$ be a domain. Recall that an $A$-module $M$ is \textbf{torsion-free}
if for every nonzero $x \in A$, the map $\mu_x: m \mapsto xm$ is injective.
\begin{enumerate}
    \item[(a)] Show that every flat $A$-module is torsion-free.
    \item[(b)] Suppose that $A$ is a PID. Show that every torsion-free
    $A$-module is flat.
    \item[(c)] Give example of a domain $A$ and torsion-free $A$-module which is not flat.
\end{enumerate}
\end{problem}

\textbf{Solution.}

\textit{Part (a).} Consider the exact sequence:
\begin{equation*}
    0 \to A \xrightarrow{\mu_x} A.
\end{equation*}

Let us apply the $(\cdot) \otimes_A M$ functor (which leaves sequence exact since
$A$ is flat):
\begin{equation*}
    0 \to M \xrightarrow{\mu_x} M.
\end{equation*}

Thus $\mu_x: M \to M, m \mapsto xm$ is injective, and so $M$ is torsion-free.\qedwhite

\textit{Part (b).} As noted in the hint, assume first that $M$ is a finitely
generated $A$-module with $A$ being a PID. Due to the Structure Theorem, $M
\cong A^r \oplus T$ with $T = \bigoplus_{i=1}^m (A/a_i)$ for $a_1 \mid \dots
\mid a_m$. In particular, when $M$ is torsion-free, $M$ is free, so $M \cong
A^r$. But free modules are always flat, so this part is finished.

Now consider the general case, where $M$ is not necessarilly finitely generated.
However, $M$ is a direct limit of finitely generated submodules: say, $M =
\varinjlim M_{\alpha}$. Each $M_{\alpha}$ is clearly torsion-free, so from the
first half, we know that each $M_{\alpha}$ is flat. Now, informally, we are left
to show that the filtered colimit of flat modules is, again, a flat module.

Consider an exact sequence $0 \to N' \to N \to N'' \to 0$. Since each
$M_{\alpha}$ is flat, we get an exact sequence $0 \to N' \otimes_A M_{\alpha} \to
N \otimes_A M_{\alpha} \to N'' \otimes_A M_{\alpha} \to 0$. Then apply the
filtered colimit to this sequence. Note two things: (1) tensor product commutes
with filtered colimits, so $N \otimes_A M \cong \varinjlim (N \otimes_A
M_{\alpha})$, (2) colimit functor is exact. Combining these two facts, we get 
an exact sequence
\begin{equation*}
    0 \to N' \otimes_A M \to N \otimes_A M \to N'' \otimes_A M \to 0,
\end{equation*}
so $(\cdot) \otimes_A M$ functor is exact, so $M$ is flat. \qedwhite

\textit{Part (c).} Take $A := k[x,y]$ and $M := (x,y)$. Clearly, $A$ is
a domain and $M$ is torsion-free. However, $M$ is not flat. Indeed, if that was
the case, for every finitely generated ideal $I \subseteq A$, the map $I
\otimes_A M \hookrightarrow M$ must be injective\footnote{\url{https://stacks.math.columbia.edu/tag/00H9}, Lemma 10.39.5(3)}. However, take $I = (x,y)$.
Then, $\gamma: (x,y) \otimes (x,y) \to (x,y), f \otimes g \mapsto fg$ is not
injective. Indeed, notice that $0 \neq x \otimes y - y \otimes x \in
\text{ker}(\gamma)$.

\begin{problem}[Exercise 4]
    Find an example of a homomorphism between domains $A \to B$ and a
torsion-free $A$-module $M$ such that $M \otimes_A B$ is not torsion-free.
\end{problem}

\textbf{Solution.} Let $A = k[x,y]$ and $B = k[x]$, where the homomorphism
$k[x,y] \to k[x]$ is given by $f(x,y) \mapsto f(x,0)$. Let $M = (x,y)$. Let us
show that $M \otimes_A B$ is not torsion-free. We have:
\begin{align*}
M \otimes_A B &= (x,y) \otimes_{k[x,y]} k[x] \\
&\cong (x, y) \otimes_{k[x,y]} (k[x,y] / (y)) \\
&\cong (x, y) / ((x,y) \cdot (y)) \\
&\cong (x,y) / (xy, y^2)
\end{align*}

However, this quotient is clearly not torsion-free. Indeed, take the element $z
:= y \in M \otimes_A B$ and $x \in B$. Then, $zx = 0$ in $M \otimes_A B$, so the
multiplication-by-$x$ map is not injective in $M \otimes_A B$.

\begin{problem}[Exercise 5]
    This problem consists of two parts:
    \begin{enumerate}
        \item[(a)] Let $I$ be a finitely generated ideal in a local ring
        $(A,\mathfrak{m})$. Prove that if $I = I^2$, then $I=A$ or $I=0$.
        \item[(b)] Let $A = C(\mathbb{R}, \mathbb{R})$ be the ring of 
        continuous functions $\mathbb{R} \to \mathbb{R}$ and $\mathfrak{m} = \{f \in A: f(0)=0\}$.
        Prove that $\mathfrak{m} = \mathfrak{m}^2$ and conclude that the ideal $\mathfrak{m}$
        is not finitely generated.
    \end{enumerate}
\end{problem}

\textbf{Solution.}

\textit{Part (a).} Suppose $I \subseteq \mathfrak{m}$ (otherwise $I=A$, and so
we are done). Then, $I = I^2 \subseteq \mathfrak{m}I$. Apply Nakayama's lemma to
get $I=0$. 

Alternatively, we can show this directly: suppose $I = (x_1,\dots,x_n)$. Since
$I \subseteq \mathfrak{m}I$, $x_i = \sum_{j=1}^n m_{ij}x_j$ for $m_{ij} \in
\mathfrak{m}$. Then, we have $\mathbf{x} = M\mathbf{x}$ for $\mathbf{x} \in A^n$
consisting of $x_1,\dots,x_n$. Note that the matrix $T := E_{n \times n}-M$ is
invertible. Indeed, consider the natural projection map $\pi: A \to
A/\mathfrak{m}$, and so $\pi(T) = \pi(E_{n \times n}) = E_{n \times n}$. Thus,
$\det(T) = 1$ in $A/\mathfrak{m}$, so in particular the determinant is not in
$\mathfrak{m}$, and so $\det(T)$ is a unit in $A$. This implies that $T$ is
invertible, so the equation $T\mathbf{x}=0$ has the only solution $\mathbf{x} =
0$. Thus, $I=0$. \qedwhite

\textit{Part (b).} Clearly, $\mathfrak{m}^2 \subseteq \mathfrak{m}$, since $I^2
\subseteq I$ holds for any ideal $I$. For $\mathfrak{m} \subseteq
\mathfrak{m}^2$, take $f \in C(\mathbb{R},\mathbb{R})$ with $f(0)=0$.
Choose
\begin{equation*}
    g(x) := \text{sign}(f(x))\sqrt{|f(x)|}, \quad h(x) := \sqrt{|f(x)|}
\end{equation*}

Notice that $g \in \mathfrak{m}$ (this function is continuous and $g(0)=0$), and
similarly $h \in \mathfrak{m}$. Moreover, $f=gh$. Thus
$\mathfrak{m}^2=\mathfrak{m}$. 

Now suppose $\mathfrak{m}$ is finitely generated. Consider the localization
$A_{\mathfrak{m}}$. Since localization preserves finite generation, either
$\mathfrak{m}A_{\mathfrak{m}}=0$ or
$\mathfrak{m}A_{\mathfrak{m}}=A_{\mathfrak{m}}$. The latter is impossible, so
$\mathfrak{m}A_{\mathfrak{m}}=0$. This means that there exists some continuous
$f$ such that $f(0) \neq 0$ and $f\cdot g = 0$ for every $g \in \mathfrak{m}$.
This is clearly impossible. Denote by $B_{\varepsilon}(x_0) = \{x \in
\mathbb{R}: |x-x_0|<\varepsilon\}$. Then,
\begin{enumerate}
    \item[(1)] there exists a neighbourhood $B_{\varepsilon}(0)$ such that
$B_{\varepsilon}(0) \subset \text{supp}(f)$;
    \item[(2)] we can pick $g \in A$ such that $\text{supp}(g) =
B_{|x_0|/2}(x_0) \not\ni 0$ for $x_0 \in B_{\varepsilon/2}(0)$\footnote{For
instance, take a bump function:
\url{https://ncatlab.org/nlab/show/bump+function}.}.
\end{enumerate}

Then, $g(0)=0$, so $g \in \mathfrak{m}$. Moreover, $f(x_0)g(x_0) \neq 0$, and so
$fg \neq 0$. \qedwhite

\end{document}